\chapter{Of the Cleaning}
\label{art:cleaning}

\articlenote{In which is set down the gravest sentence this Pact names: the sending of a citizen outside to clean, that the silo may endure and every other citizen may live.}

\section{The Request and the Right of a Citizen}

No citizen may be compelled to walk outside the silo, to touch the lenses or sensors, or to approach the airlock against that citizen's will. Yet any citizen, at any time, may request to go outside, and that request must be granted. This choice, once made, is final and irrevocable. A citizen granted this request is released from all bonds of the silo and all claims to the rolls, and is known thereafter as a cleaner.

\begin{clauses}
\item A citizen may make a final request at any point in the citizen's life, to any officer of Judicial or to the Sheriff, or may write it down and place it in the archive under Article 6, by speaking or setting down these words and no others: ``I wish to walk outside.'' No further explanation is required of the citizen, and none shall be asked; the words alone, once said in earnest and to a proper officer, are the whole of the request this Pact requires.
\item The request does not need to be approved, discussed, or delayed; it shall be granted with utmost expedience.
\item A citizen who makes a final request shall be brought to the airlock on the next available cycle, clothed in a suit provided by Information Technology under Section 3 of this Article, and the citizen shall walk outside. No Pact provides a method by which a citizen might return, and none is needed, for every request is final.
\item A citizen who walks outside remains a citizen in the rolls until Judicial is informed by Information Technology under Article 8 that the sensors confirm the citizen's death; only then is the citizen's name closed under Article 2.
\end{clauses}

\section{The Obligation to Clean and the Sentencing to Cleaning}

A citizen sentenced to cleaning is not condemned as a murderer is condemned; rather, a citizen sentenced to cleaning is a citizen who has violated the Pact in a way that threatens the whole silo, and the silo has determined that such a citizen is a danger to all others and must be removed. This is not a death sentence; it is an exile.

\begin{clauses}
\item The Judge alone sentences a citizen to cleaning, for an offense named among the gravest under Article 17. No citizen sentenced to cleaning may be reprieved, pardoned, or offered mercy by any other office.
\item A citizen sentenced to cleaning shall be prepared by the Infirmary under Article 19 for the journey outside: clothed in a suit, provided such water or food as that preparation requires, and made ready without delay.
\item A citizen sentenced to cleaning is brought to the airlock and is sent outside on the next available cycle, exactly as a citizen making a final request, and the two are thereafter indistinguishable: both are cleaners, and both are lost to the silo.
\end{clauses}

\section{The Suits and the Seals}

The Pact does not name the fate of a citizen who walks outside, nor does it describe how long a citizen might last in the air above. The silo's knowing or not knowing of that fate is no citizen's concern. The suit provided is assembled in the cleaning lab under the charge of Information Technology, from cloth, seals, and fittings supplied by Mechanical and Supply, tested to hold for such time as the silo judges reasonable, and made ready with the utmost care.

\begin{clauses}
\item The suit provided to any citizen sent to clean is assembled and sealed to a specification set by Information Technology, exercised under such further procedure as that Department alone is instructed in, this Pact not setting down that procedure here, for it is no part of ordinary citizenship. Mechanical and Supply furnish the cloth, seals, and fittings the specification requires, but the specification itself, and the final sealing, are Information Technology's alone. A suit known to be faulty shall not be sent out.
\item Information Technology under Article 8 monitors the sensors after a citizen emerges to confirm the citizen has begun and completed the cleaning, and reports this finding to Judicial for the archive under Article 6. The citizen is not informed of the manner of their death, nor of the time it takes.
\item A citizen being prepared for cleaning retains the right to refuse the suit; such a refusal is not grounds for a stay or a reprieve, and the citizen is brought outside without the suit.
\item The citizens of Information Technology who assemble, seal, and administer the suit and the chamber are entered into that trade only by full shadowing under Article 13, and are vetted beyond the ordinary measure Article 13 there provides, as befits a trade nearer than any other to the silo's gravest and most final procedure. The particulars of that further vetting are kept by the Head of Information Technology and are no part of ordinary citizenship.
\end{clauses}

\section{The Airlock and the Departure}

The airlock stands at the silo's uppermost level, beyond the wallscreen and its cafeteria, and it is through this chamber alone that a citizen sent to clean departs the silo. This section sets down the plain administration of that departure, that it be carried out with the same order the silo asks of every other procedure it keeps.

\begin{clauses}
\item Before the outer door is opened, the citizen is given no fewer than three and no more than five pads of raw wool, drawn from a store kept for this purpose alone and logged, by number issued and number remaining, in the chamber's own record, that the citizen may, once outside, wipe clean the lenses of the outward sensors, so that the wallscreen's image may continue to be gathered as true as it can be kept. A citizen sent to clean under sentence performs this wiping before any other act; a citizen departed by final request performs it if the citizen is willing, and is asked no further than once.
\item Before the flush begins, the Officer of Judicial or the Sheriff who brought the citizen to the chamber confirms the inner door is sealed, enters the hour of sealing in the chamber's record by the silo's own clock, and reports the sealing to the technician of Information Technology who oversees the chamber, who countersigns the same entry. The chamber is then sealed and flushed of the silo's own air, by a gas Information Technology alone administers and this Pact does not further describe, that the poison of the world above not be carried back into the silo with the door's opening; the technician enters the flush's completion in the same record, by the same hour and countersignature this section already requires.
\item Only once the technician of Information Technology has made the entry required above, confirming the chamber is flushed and the citizen is suited and sealed, does the outer door open. The Officer or Sheriff who accompanied the citizen does not enter the chamber for this final step, and returns to ordinary duty once the inner door is sealed behind the citizen and the record's first entry is made.
\end{clauses}

\section{The Public Witnessing of a Cleaning}

A cleaning at the wallscreen of the cafeteria is neither compulsory for any citizen to attend, nor forbidden. The citizens who gather to watch do so of their own choice, and the Sheriff ensures peace among them as any other gathering under Article 7.

\begin{clauses}
\item A cleaning is displayed on the wallscreen at the Level 1 cafeteria as it occurs, exactly as the camera and the sensors capture it, without interruption or interference, so that all who wish to know may see the fate of the citizen sent out.
\item The gathering to watch is the charge of the Sheriff under Article 7, and the Sheriff shall keep order as at any other public place, and may disperse a crowd that threatens to become disorder, though the gathering itself may not be forbidden.
\item After a cleaning is complete, the image on the wallscreen ceases, and the visual record of what transpired outside is kept in the archive under Article 6, preserved for the record but not displayed further.
\end{clauses}

\section{The Ritual and the Meaning}

The cleaning is not cruelty, though it appears so to the citizen sent to it. It is the silo's answer to those offenses that endanger not a single citizen but all citizens at once. The sensor tampering, the interference with vital systems, the self-killing of a citizen, the forging of the rolls --- these strike at the whole, and the whole must answer.

\begin{clauses}
\item No citizen shall be sent to cleaning for a lesser crime than those named under Article 17, Section 4. No Judge may increase this severity, nor may any other office. This severity is fixed and irrevocable.
\item The families of citizens sent to cleaning are not barred from continuing to live in the silo, nor marked, nor punished under this Pact, save where the family member themselves has violated the Pact. A child born to a parent who was cleaned is entered in the rolls exactly as any other child.
\item A citizen watching a cleaning, or witnessing its aftermath in the archive, shall draw such meaning as the citizen finds there. This Pact provides the mechanism but not the interpretation. The silo persists because such measures are taken, and every citizen's life depends on this.
\end{clauses}
